\chapter{The Classical-to-Quantum Bridge}

This final chapter looks forward from the classical geometric framework we have built toward the quantisation of the electromagnetic field. We do not develop full quantum electrodynamics --- that would require a separate volume --- but we identify exactly which mathematical structures from the classical theory survive, which must be modified, and which new structures emerge. The goal is to ensure that the reader knows precisely what is needed to ``do QED honestly'' and where each piece of the classical bundle apparatus reappears in the quantum theory.

\section{The Classical Configuration Space}

\subsection{The space of connections}

In the classical theory, the dynamical variable is the connection $A$ on a fixed principal $U(1)$-bundle $P \to M$. The space of all connections is
\begin{equation}
\mathcal{A} = \{\text{connections on } P\}.
\end{equation}

As noted in Chapter~7, $\mathcal{A}$ is an infinite-dimensional \textbf{affine space} modelled on $\Omega^1(M, \Lie{u}(1)) \cong \Omega^1(M)$. It is the classical analogue of the phase space of a finite-dimensional mechanical system.

\subsection{The gauge group}

The group of gauge transformations is
\begin{equation}
\mathcal{G} = \Ck(M, U(1)) = \{\text{smooth maps } M \to U(1)\}.
\end{equation}

$\mathcal{G}$ acts on $\mathcal{A}$ by $A \mapsto A + \dd\chi$ (for $g = e^{i\chi} \in \mathcal{G}$). Physically equivalent configurations are related by gauge transformations, so the true configuration space is the \textbf{orbit space}:
\begin{equation}
\mathcal{B} = \mathcal{A} / \mathcal{G}.
\end{equation}

\begin{remark}
The orbit space $\mathcal{B}$ is generally \emph{not} a smooth manifold --- it may have singularities at connections with non-trivial stabilisers (``reducible connections''). For $U(1)$, the stabiliser is always $U(1)$ itself (constant gauge transformations), so $\mathcal{B}$ has a mild orbifold structure. For non-abelian groups, the singularity structure is richer and physically significant.
\end{remark}

\section{The Path Integral over Connections}

\subsection{The Feynman path integral}

The quantum theory is (heuristically) defined by the \textbf{path integral}:
\begin{equation}
Z = \int_{\mathcal{A}/\mathcal{G}} \mathcal{D}[A]\; e^{iS[A]/\hbar},
\end{equation}
where $S[A] = -\frac{1}{8\pi}\int_M F \wedge \hodge F$ is the Yang--Mills action, and $\mathcal{D}[A]$ denotes a (formal) measure on the space of gauge-equivalence classes of connections.

The path integral sums over \emph{all} connections modulo gauge, weighted by the classical action. Classical solutions (those extremising $S$) dominate in the $\hbar \to 0$ limit; quantum fluctuations around them contribute corrections.

\subsection{Gauge fixing}

In practice, one does not integrate over $\mathcal{A}/\mathcal{G}$ directly. Instead, one chooses a \textbf{gauge-fixing condition} (e.g.\ Lorenz gauge $\partial_\mu A^\mu = 0$, Coulomb gauge $\vec{\nabla}\cdot\vec{A} = 0$) that selects a unique representative from each gauge orbit. The \textbf{Faddeev--Popov procedure} implements this by introducing ghost fields and a determinant factor that corrects for the redundancy:
\begin{equation}
Z = \int_{\mathcal{A}} \mathcal{D}A\;\delta(\text{gauge condition})\;\det\!\left(\frac{\delta(\text{gauge condition})}{\delta\chi}\right)\;e^{iS[A]/\hbar}.
\end{equation}

For $U(1)$, the Faddeev--Popov determinant is a constant (independent of $A$), and ghosts decouple. This is another simplification special to the abelian case; for non-abelian gauge theories, the ghosts are dynamical and essential.

\section{Topology of the Configuration Space and Theta-Sectors}

\subsection{Large gauge transformations}

The gauge group $\mathcal{G} = \Ck(M, U(1))$ has a connected component of the identity $\mathcal{G}_0$ (consisting of gauge transformations continuously connected to the identity) and possibly disconnected components:
\begin{equation}
\pi_0(\mathcal{G}) = \mathcal{G}/\mathcal{G}_0.
\end{equation}

A gauge transformation in $\mathcal{G} \setminus \mathcal{G}_0$ is called a \textbf{large gauge transformation}. It cannot be continuously deformed to the identity.

\begin{example}
On $M = S^1 \times \R^3$ (thermal field theory), $\mathcal{G} = \Ck(S^1 \times \R^3, U(1))$ and large gauge transformations are classified by the winding number of the map $S^1 \to U(1)$: $\pi_0(\mathcal{G}) \cong \pi_1(U(1)) \cong \Z$.
\end{example}

\subsection{Theta-vacua}

If $\pi_0(\mathcal{G}) \neq 0$, the configuration space $\mathcal{B} = \mathcal{A}/\mathcal{G}$ is not simply connected. The path integral must account for the topological sectors:
\begin{equation}
Z(\theta) = \sum_{n \in \pi_0(\mathcal{G})} e^{in\theta} \int_{\mathcal{A}_n/\mathcal{G}_0} \mathcal{D}[A]\;e^{iS[A]/\hbar},
\end{equation}
where $\theta \in [0, 2\pi)$ is the \textbf{vacuum angle} and $\mathcal{A}_n$ denotes the sector with topological charge $n$.

For $U(1)$ in four-dimensional Minkowski spacetime, $\pi_0(\mathcal{G}) = 0$ and there is no theta-parameter. The theta-vacuum structure becomes non-trivial for:
\begin{itemize}
\item Non-abelian gauge theories ($SU(2)$, $SU(3)$): $\pi_0(\mathcal{G}) \cong \pi_3(G) \cong \Z$, leading to the \textbf{strong CP problem} in QCD.
\item $U(1)$ on non-trivial spacetimes: if $M$ has non-trivial topology, large gauge transformations and theta-sectors can appear even for the abelian theory.
\end{itemize}

\section{The Geometric Structure that Survives Quantisation}

\begin{table}[h]
\centering
\renewcommand{\arraystretch}{1.5}
\begin{tabular}{lll}
\toprule
\textbf{Classical} & \textbf{Quantum} & \textbf{Status} \\
\midrule
Principal $U(1)$-bundle $P$ & Same & Unchanged \\
Connection $\omega$ & Path integral over connections & Promoted to quantum variable \\
Curvature $F = \dd A$ & Field operator $\hat{F}$ & Operator-valued 2-form \\
Gauge invariance & Ward--Takahashi identities & Exact quantum symmetry \\
Bianchi identity $\dd F = 0$ & Operator identity & Unchanged (kinematic) \\
$\dd\,\hodge F = \frac{4\pi}{c}\hodge J$ & Heisenberg equation of motion & Quantum dynamics \\
Holonomy $\exp(i\oint A)$ & Wilson loop operator & Fundamental observable in LGT \\
Chern number $c_1 \in \Z$ & Topological charge & Labels superselection sectors \\
Associated bundle $E_n$ & Fock space of charge-$n$ particles & Quantum matter \\
Covariant derivative $\nabla$ & Interaction vertex $e\bar{\psi}\gamma^\mu A_\mu\psi$ & QED coupling \\
\bottomrule
\end{tabular}
\end{table}

Every element of the classical geometric framework reappears in the quantum theory, often in a more sophisticated guise. The bundle structure, the connection, the curvature, gauge invariance, and characteristic classes all survive quantisation. What changes is the \emph{dynamics}: classical trajectories (extrema of $S$) are replaced by the sum over all configurations (the path integral).

\section{Anomalies: When Classical Symmetries Fail}

A \textbf{quantum anomaly} occurs when a symmetry of the classical action is not preserved by the quantisation procedure (specifically, by the path integral measure $\mathcal{D}A$).

\begin{definition}[Gauge anomaly]
A \textbf{gauge anomaly} in a theory with gauge group $G$ is the failure of gauge invariance at the quantum level:
\begin{equation}
\delta_\chi Z \neq 0.
\end{equation}
\end{definition}

For $U(1)$ electrodynamics, the gauge anomaly cancels provided the sum of the charges of all fermion species vanishes: $\sum_f n_f = 0$. In the Standard Model, this condition is satisfied within each generation of quarks and leptons --- a non-trivial constraint that the particle content must satisfy for the theory to be consistent.

The anomaly is computed from the \textbf{index theorem} (Atiyah--Singer), which relates the analytical properties of the Dirac operator $\slashed{D} = \gamma^\mu\nabla_\mu$ to the topology of the bundle:
\begin{equation}
\text{index}(\slashed{D}) = \int_M \hat{A}(M) \wedge \mathrm{ch}(E),
\end{equation}
where $\hat{A}(M)$ is a characteristic class of the tangent bundle (related to gravity) and $\mathrm{ch}(E)$ is the Chern character of the associated bundle (related to the gauge field). The interplay between gravitational and gauge topology is a deep feature of the quantum theory that is invisible classically.

\section{From Connections to Photons}

The quantisation of the free electromagnetic field (in Lorenz gauge, on Minkowski spacetime) proceeds by expanding $A$ in Fourier modes and promoting the coefficients to creation and annihilation operators. The resulting Fock space describes states with definite numbers of \textbf{photons} --- the quanta of the electromagnetic field.

The two physical polarisation states of the photon correspond to the two transverse degrees of freedom of the connection, after gauge-fixing removes the longitudinal and temporal modes. The gauge invariance of the classical theory (the redundancy of $A \to A + \dd\chi$) ensures that the unphysical polarisations decouple from all physical amplitudes --- this is the content of the Ward--Takahashi identity.

\section{The Geometric Unity of the Standard Model}

The full Standard Model of particle physics is a Yang--Mills gauge theory with structure group
\begin{equation}
G_{\text{SM}} = SU(3) \times SU(2) \times U(1).
\end{equation}

Every piece of the geometric framework we have developed for $U(1)$ electrodynamics generalises directly:

\begin{itemize}
\item The gauge field is a connection on a principal $G_{\text{SM}}$-bundle over spacetime.
\item The field strength is the curvature of this connection.
\item Quarks and leptons are sections of associated vector bundles determined by their representations under $G_{\text{SM}}$.
\item The covariant derivative encodes all gauge interactions (electromagnetic, weak, strong).
\item The action is the Yang--Mills functional $S = \int \tr(F \wedge \hodge F)$ plus matter terms.
\item Topological effects (instantons, theta-vacua, anomaly cancellation) are governed by the characteristic classes of the $G_{\text{SM}}$-bundle.
\end{itemize}

The language of principal bundles, connections, and curvature is not merely a convenient reformulation of electromagnetism --- it is the mathematical framework of all known fundamental interactions except gravity. (And even gravity, in the Cartan formulation, is a connection on a frame bundle.)

\section{Conclusion}

We began this book with bare sets and topological spaces. We built smooth manifolds, populated them with differential forms, and showed that half of Maxwell's equations ($\dd F = 0$) are purely topological --- they hold without any metric. We introduced the metric and the Hodge star to write the dynamical half ($\dd\,\hodge F = \hodge J$). We then went deeper, recognising the gauge potential as a connection on a principal $U(1)$-bundle, the field strength as its curvature, gauge invariance as the freedom of choosing local sections, and charge quantisation as a topological theorem about the classification of bundles.

The structuralist principle --- \emph{never introduce a physical concept without first having the mathematical object that carries it} --- has revealed that electromagnetism is not a collection of empirical laws but a transparent branch of differential geometry. Every computational tool of classical electrodynamics (potentials, gauge fixing, the Lorentz force, the energy--momentum tensor) has a precise home in this geometric hierarchy, and the distinction between structure and convention is always clear.

The geometric language does not merely reformulate known results in fancier notation. It \emph{explains} phenomena --- the Aharonov--Bohm effect, charge quantisation, the impossibility of a global potential for monopoles --- that are mysterious or ad hoc in the traditional formulation. And it provides the exact mathematical tools needed to extend electromagnetism to non-abelian gauge theories, to curved spacetime, and to the quantum domain.

This is the geometrical anatomy of electrodynamics.
